\section{Higgs 真空与电弱破缺}\label{sec:sm-higgs-ewsb-dp11}

手征规范结构在未破缺真空中给出一个\emph{无质量的规范理论}：规范对称性允许 Yukawa 三场算符，却仍禁止直接写 $m_W^2W_\mu W^\mu$、$m_Z^2Z_\mu Z^\mu$ 和费米子裸 Dirac 质量。这里的\emph{真空}不是“所有场都等于零”的同义词，而是量子场论中最低能量的参考态；若某个场算符在这个最低能态中的期望值非零，就称它具有非零真空期望值（vacuum expectation value, VEV）。标准模型的关键正是 Higgs 场的最低能构型具有非零模长。

拉格朗日量仍具有完整 $SU(2)_L\times U(1)_Y$ 对称性，但选定一个具体最低能真空以后，只有把这个真空保持不变的群变换仍显式作用在真空上。这个“作用量对称、真空只保持子群”的结构称为自发对称性破缺。普通力学中可把它类比成一圈等高的连续极小值：方程对绕中心旋转对称，但任选一个极小点作为展开中心后，该点本身只被较小的稳定子保持。对标准模型，先求 Higgs 势的真空流形，再把 $|D_\mu\Phi|^2$ 在真空上展开。势能沿真空流形切向没有恢复力的方向对应 Goldstone 坐标；在规范理论中，这三个坐标与三个有质量规范场的纵向自由度重组，而径向涨落保留下来成为物理 Higgs 场。对角化规范场质量矩阵以后，唯一的零本征方向才被识别为电磁 $U(1)$，此时电荷算符 $Q$ 和弱混合角才有明确物理意义。

\subsection{Higgs 真空与径向模}

Higgs 二重态包含四个实场自由度。势能先决定真空模长；选定一个真空方向以后，沿真空轨道切向的三个方向对应被破缺生成元，径向方向则是具有恢复力的物理标量涨落。随后把同一个真空代入协变动能，三个切向自由度与三个有质量规范场的纵向极化结合，而径向模保留下来成为物理 Higgs 场。

对 \eqnrefs{eq:sm-higgs-rep-dp11} 的复二重态 $\Phi$，洛伦兹不变性、$SU(2)_L\times U(1)_Y$ 规范不变性与可重整化性把标量拉氏量限制为
\begin{equation}
 \boxed{
 \Lag_{\rm Higgs}
 =(D_\mu\Phi)^\dagger(D^\mu\Phi)
 -V(\Phi),
 \qquad
 V(\Phi)=-\mu_H^2\Phi^\dagger\Phi
 +\frac{\lambda_H}{\hbar c}(\Phi^\dagger\Phi)^2.}
 \label{eq:sm-higgs-lagrangian-dp11}
\end{equation}
取 $\mu_H^2>0$、$\lambda_H>0$。由于势能只依赖规范不变量 $r^2\equiv\Phi^\dagger\Phi$，真空首先由这个径向变量的极小值决定。非零极小值满足 $r^2=\mu_H^2\hbar c/(2\lambda_H)$，因此可把真空代表元与能标定义为
\begin{equation}
 \boxed{
 \langle\Phi\rangle
 =\frac1{\sqrt2}\begin{pmatrix}0\\v\end{pmatrix},
 \qquad
 v^2=\frac{\mu_H^2\hbar c}{\lambda_H},
 \qquad
 v_E\equiv\sqrt{\hbar c}\,v.}
 \label{eq:sm-higgs-vev-dp11}
\end{equation}
$v_E$ 具有能量单位，是电弱质量与 Higgs 耦合共同使用的真空能标。

一般场构型可沿规范轨道与径向方向分解为
\begin{equation*}
 \Phi(x)
 =\frac1{\sqrt2}
 U_\xi(x)
 \begin{pmatrix}0\\v+h(x)\end{pmatrix},
 \qquad
 U_\xi(x)=\exp\!\left[
 \frac{\ii\xi^I(x)\sigma^I}{v}
 \right].
\end{equation*}
$\xi^I$ 描述三个规范轨道切向自由度。所谓\emph{幺正规范}（unitary gauge）是利用局域 $SU(2)_L\times U(1)_Y$ 规范自由度，把这三个沿规范轨道的坐标选为零，即取 $U_\xi=I$。在这种规范下，被破缺方向的 Goldstone 场不再显式出现在标量二重态中，它们对应的自由度已经进入 $W^\pm$ 与 $Z$ 的纵向极化；标量扇区只剩径向场 $h$：
\begin{equation}
 \Phi_U(x)=\frac1{\sqrt2}\begin{pmatrix}0\\v+h(x)\end{pmatrix}.
 \label{eq:sm-higgs-unitary-gauge-dp11}
\end{equation}
真空附近的质量与自相互作用由势能沿径向方向的展开系数决定：
\begin{equation}
 \boxed{
 V(h)=V_0+\mu_H^2h^2
 +\frac{\lambda_Hv}{\hbar c}h^3
 +\frac{\lambda_H}{4\hbar c}h^4.}
 \label{eq:sm-higgs-potential-expanded-dp11}
\end{equation}
二次项的曲率定义物理 Higgs 质量，得到
\begin{equation}
 \boxed{
 m_hc^2=\sqrt{2\lambda_H}\,v_E.}
 \label{eq:sm-higgs-mass-dp11}
\end{equation}
因此 Higgs 质量固定四次耦合常数：
\begin{equation*}
 \lambda_H=\frac{(m_hc^2)^2}{2v_E^2}.
\end{equation*}

\begin{figure}[htbp]
 \centering
 \includegraphics[width=.72\textwidth]{sm_higgs_potential_dp11.png}
 \caption{Higgs 势能的径向截面。$\mu_H^2>0$ 时原点是不稳定驻点，连续极小值满足 $\Phi^\dagger\Phi=v^2/2$。规范固定选定一个真空代表元后，径向曲率给出物理 Higgs 质量；规范轨道上的三个切向自由度与有质量弱规范玻色子的纵向偏振重组。}
 \label{fig:sm-higgs-potential-dp11}
\end{figure}

\begin{derivation}[从\eqnrefs{eq:sm-higgs-lagrangian-dp11}到\eqnrefs{eq:sm-higgs-mass-dp11}]
把正文 Higgs 势写成单一规范不变量
\begin{equation*}
 x\equiv\Phi^\dagger\Phi\ge0,
 \qquad
 V(x)=-\mu_H^2x+\frac{\lambda_H}{\hbar c}x^2.
\end{equation*}
对复二重态作任意变分时，$\delta x=\delta\Phi^\dagger\Phi+\Phi^\dagger\delta\Phi$，因此驻点满足
\begin{equation*}
 \frac{\partial V}{\partial\Phi^\dagger}
 =\left(-\mu_H^2+\frac{2\lambda_H}{\hbar c}\Phi^\dagger\Phi\right)\Phi=0.
\end{equation*}
除原点 $\Phi=0$ 外，非零驻点组成
\begin{equation*}
 \Phi^\dagger\Phi
 =\frac{\mu_H^2\hbar c}{2\lambda_H}.
\end{equation*}
当 $\mu_H^2>0$、$\lambda_H>0$ 时，原点沿径向的二阶变化为负，而上述非零流形是稳定极小值。利用电弱规范变换可以从这条真空轨道中选取代表元
\begin{equation*}
 \langle\Phi\rangle
 =\frac1{\sqrt2}\binom0v,
\end{equation*}
于是
\begin{equation*}
 \frac{v^2}{2}
 =\frac{\mu_H^2\hbar c}{2\lambda_H},
 \qquad
 v^2=\frac{\mu_H^2\hbar c}{\lambda_H},
 \qquad
 v_E\equiv\sqrt{\hbar c}\,v,
\end{equation*}
即正文\eqnrefs{eq:sm-higgs-vev-dp11}。

真空附近先保留一般规范轨道方向，写成
\begin{equation*}
 \Phi(x)=\frac1{\sqrt2}U_\xi(x)\binom0{v+h(x)}.
\end{equation*}
势能只依赖 $\Phi^\dagger\Phi$，所以 $U_\xi^\dagger U_\xi=I$ 使三个切向场 $\xi^I$ 完全不进入 $V$；势能只约束径向变量 $h$。取幺正规范 $U_\xi=I$ 后，
\begin{equation*}
 \Phi^\dagger\Phi=\frac12(v+h)^2,
\end{equation*}
代回势能：
\begin{align*}
 V(h)
 &=-\frac{\mu_H^2}{2}(v+h)^2
 +\frac{\lambda_H}{4\hbar c}(v+h)^4\\
 &=-\frac{\mu_H^2}{2}
 (v^2+2vh+h^2)
 +\frac{\lambda_H}{4\hbar c}
 (v^4+4v^3h+6v^2h^2+4vh^3+h^4).
\end{align*}
逐阶收集：
\begin{align*}
 V(h)=V_0
 &+v\left(-\mu_H^2+\frac{\lambda_Hv^2}{\hbar c}\right)h\\
 &+\left(-\frac{\mu_H^2}{2}+\frac{3\lambda_Hv^2}{2\hbar c}\right)h^2
 +\frac{\lambda_Hv}{\hbar c}h^3
 +\frac{\lambda_H}{4\hbar c}h^4.
\end{align*}
利用真空条件 $\lambda_Hv^2/(\hbar c)=\mu_H^2$，线性项消失，二次项化为
\begin{equation*}
 V(h)=V_0+\mu_H^2h^2
 +\frac{\lambda_Hv}{\hbar c}h^3
 +\frac{\lambda_H}{4\hbar c}h^4,
\end{equation*}
即正文\eqnrefs{eq:sm-higgs-potential-expanded-dp11}。

最后核对径向场的动能归一化。忽略规范场时，
\begin{equation*}
 (\partial_\mu\Phi_U)^\dagger(\partial^\mu\Phi_U)
 =\frac12\partial_\mu h\,\partial^\mu h,
\end{equation*}
所以 $h$ 已经是标准规范化的实标量。标准质量项为
\begin{equation*}
 -\frac12\left(\frac{m_hc}{\hbar}\right)^2h^2.
\end{equation*}
由于拉格朗日量含 $-V$，比较二次项得到
\begin{equation*}
 \frac12\left(\frac{m_hc}{\hbar}\right)^2=\mu_H^2.
\end{equation*}
再用
\begin{equation*}
 \mu_H^2=\frac{\lambda_Hv^2}{\hbar c}
 =\frac{\lambda_Hv_E^2}{(\hbar c)^2},
\end{equation*}
便有
\begin{equation*}
 (m_hc^2)^2=2\lambda_Hv_E^2,
 \qquad
 \boxed{m_hc^2=\sqrt{2\lambda_H}\,v_E},
\end{equation*}
即正文\eqnrefs{eq:sm-higgs-mass-dp11}。同一展开同时固定 $h^3$ 和 $h^4$ 自耦，因此质量、三点自耦和四点自耦都来自同一个径向势，而不是彼此独立的参数。
\end{derivation}

\subsection{电弱质量矩阵}

Higgs 场属于颜色表示 $\mathbf1$，所以 $SU(3)_c$ 生成元在该表示上作用为零；协变导数中的颜色联络项因此消失，留下电弱规范场：
\begin{equation*}
 D_\mu\Phi
 =\left[
 \partial_\mu
 +\ii g\frac{\sigma^I}{2}\mathcal W_\mu^I
 +\ii g'\frac12\mathcal B_\mu
 \right]\Phi.
\end{equation*}
把真空\eqnrefs{eq:sm-higgs-vev-dp11} 代入 $|D_\mu\Phi|^2$ 后，规范场的二次型在对角化之前统一写成
\begin{align}
 (D_\mu\langle\Phi\rangle)^\dagger(D^\mu\langle\Phi\rangle)
 =&\frac{v^2g^2}{8}\left[(\mathcal W^1)^2+(\mathcal W^2)^2\right]\\
 &+\frac{v^2}{8}(g\mathcal W^3-g'\mathcal B)^2.
 \label{eq:sm-gauge-mass-before-diagonal-dp11}
\end{align}
这个二次型自然分成两个线性代数问题：$(\mathcal W^1,\mathcal W^2)$ 构成带电扇区，而 $(\mathcal W^3,\mathcal B)$ 构成需要对角化的 $2\times2$ 中性质量矩阵。

带电扇区由下列线性组合对角化：
\begin{equation}
 \boxed{
 \mathcal W_\mu^\pm
 =\frac{\mathcal W_\mu^1\mp\ii\mathcal W_\mu^2}{\sqrt2}.}
 \label{eq:wpm-definition-dp11}
\end{equation}
因为 $(\mathcal W^1)^2+(\mathcal W^2)^2=2\mathcal W^+\mathcal W^-$，第一行成为
\begin{equation*}
 \frac{g^2v^2}{4}\mathcal W_\mu^+\mathcal W^{-\mu}.
\end{equation*}
与几何归一化下的有质量矢量场质量项
\begin{equation*}
 \hbar c\left(\frac{m_Wc}{\hbar}\right)^2
 \mathcal W_\mu^+\mathcal W^{-\mu}
\end{equation*}
比较：
\begin{equation}
 \boxed{m_Wc^2=\frac{gv_E}{2}.}
 \label{eq:sm-w-mass-dp11}
\end{equation}

中性扇区写成矩阵形式：
\begin{equation}
 \Lag_{\rm neutral}^{\rm mass}
 =\frac{v^2}{8}
 \begin{pmatrix}\mathcal W_\mu^3&\mathcal B_\mu\end{pmatrix}
 \begin{pmatrix}g^2&-gg'\\-gg'&g'^2\end{pmatrix}
 \begin{pmatrix}\mathcal W^{3\mu}\\\mathcal B^\mu\end{pmatrix}.
 \label{eq:neutral-mass-matrix-dp11}
\end{equation}
这个矩阵的两个本征值是 $0$ 与 $g^2+g'^2$。零本征值对应未破缺的电磁方向，非零本征值对应有质量中性矢量场。为写出归一化本征矢，定义弱混合角
\begin{equation}
 \boxed{
 \sin\theta_W\equiv s_W=\frac{g'}{\sqrt{g^2+g'^2}},
 \qquad
 \cos\theta_W\equiv c_W=\frac{g}{\sqrt{g^2+g'^2}}.}
 \label{eq:weak-angle-dp11}
\end{equation}
正交变换
\begin{equation}
 \boxed{
 \begin{pmatrix}A_\mu^{\rm geo}\\\mathcal Z_\mu\end{pmatrix}
 =\begin{pmatrix}s_W&c_W\\c_W&-s_W\end{pmatrix}
 \begin{pmatrix}\mathcal W_\mu^3\\\mathcal B_\mu\end{pmatrix}}
 \label{eq:az-mixing-dp11}
\end{equation}
给出质量本征态。标准模型质量基中的光子沿用规范场的几何归一化，记为 $A_\mu^{\rm geo}$；前面 Maxwell/QED 章节的 $A_\mu$ 则是 SI 四势。二者满足
\begin{equation}
 \boxed{
 A_\mu^{\rm geo}
 =\sqrt{\frac{\varepsilon_0c}{\hbar}}\,A_\mu,
 \qquad
 F_{\mu\nu}^{\rm geo}
 =\sqrt{\frac{\varepsilon_0c}{\hbar}}\,F_{\mu\nu},
 \qquad
 g_{\rm em}A_\mu^{\rm geo}=\frac{e}{\hbar}A_\mu.}
 \label{eq:sm-photon-geo-si-bridge-dp75}
\end{equation}
其中 $F_{\mu\nu}^{\rm geo}=\partial_\mu A_\nu^{\rm geo}-\partial_\nu A_\mu^{\rm geo}$，$e>0$ 是元电荷，$g_{\rm em}=e/\sqrt{\varepsilon_0\hbar c}$。因此 $[A_\mu^{\rm geo}]={\rm m^{-1}}$，而 $A_\mu$ 保持 SI 四势量纲；\eqnrefs{eq:sm-photon-geo-si-bridge-dp75} 把标准模型无量纲耦合与 QED 的物理电荷归一化接到同一电磁场上。代回 \eqnrefs{eq:neutral-mass-matrix-dp11}：
\begin{equation}
 \boxed{
 m_\gamma=0,
 \qquad
 m_Zc^2=\frac{v_E}{2}\sqrt{g^2+g'^2}
 =\frac{m_Wc^2}{c_W}.}
 \label{eq:sm-z-photon-masses-dp11}
\end{equation}
光子无质量来自 Higgs 真空仍保留一个未破缺规范生成元：对应方向正是中性质量矩阵的零本征矢。由同一质量关系还可定义树级电弱参数
\begin{equation}
 \boxed{
 \rho_{\rm tree}
 \equiv\frac{(m_Wc^2)^2}{(m_Zc^2)^2c_W^2}=1.}
 \label{eq:sm-rho-tree-dp85}
\end{equation}
这里的等号只使用 Higgs 二重态给出的质量关系 $m_Zc^2=m_Wc^2/c_W$；圈修正或不同 Higgs 表示会使精密电弱参数偏离这一树级值。

\begin{derivation}[从\eqnrefs{eq:sm-gauge-mass-before-diagonal-dp11}到\eqnrefs{eq:sm-z-photon-masses-dp11}]
正文\eqnrefs{eq:sm-gauge-mass-before-diagonal-dp11} 已给出 Higgs 真空诱导的带电与中性规范场二次型。带电扇区直接与有质量矢量场的标准质量项比较；中性扇区则对角化 $2\times2$ 质量矩阵。两步合在一起给出正文\eqnrefs{eq:sm-z-photon-masses-dp11} 中的零质量光子以及 $W$、$Z$ 质量。

对 Higgs 真空，
\begin{equation*}
 D_\mu\langle\Phi\rangle
 =\frac{\ii v}{2\sqrt2}
 \left[g\sigma^I\mathcal W_\mu^I+g'\mathcal B_\mu\right]
 \begin{pmatrix}0\\1\end{pmatrix}.
\end{equation*}
Pauli 矩阵作用在下分量上给出
\begin{equation*}
 \sigma^1\binom01=\binom10,
 \qquad
 \sigma^2\binom01=\binom{-\ii}0,
 \qquad
 \sigma^3\binom01=\binom0{-1}.
\end{equation*}
于是
\begin{equation*}
 D_\mu\langle\Phi\rangle
 =\frac{\ii v}{2\sqrt2}
 \begin{pmatrix}
 g(\mathcal W_\mu^1-\ii\mathcal W_\mu^2)\\
 -g\mathcal W_\mu^3+g'\mathcal B_\mu
 \end{pmatrix},
\end{equation*}
取模平方即得到\eqnrefs{eq:sm-gauge-mass-before-diagonal-dp11}。

带电扇区使用 $\mathcal W^\pm=(\mathcal W^1\mp\ii\mathcal W^2)/\sqrt2$，所以
\begin{equation*}
 \frac{v^2g^2}{8}\left[(\mathcal W^1)^2+(\mathcal W^2)^2\right]
 =\frac{g^2v^2}{4}\mathcal W^+_\mu\mathcal W^{-\mu},
\end{equation*}
与有质量矢量场的质量项比较得到 $m_Wc^2=gv_E/2$。

中性扇区的无量纲矩阵为
\begin{equation*}
 M_0^2=\begin{pmatrix}g^2&-gg'\\-gg'&g'^2\end{pmatrix}.
\end{equation*}
其特征多项式为
\begin{align*}
 \det(M_0^2-\lambda I)
 &=(g^2-\lambda)(g'^2-\lambda)-g^2g'^2\\
 &=\lambda\left[\lambda-(g^2+g'^2)\right].
\end{align*}
因此本征值是 $0$ 与 $g^2+g'^2$；相应的归一化本征矢正是\eqnrefs{eq:az-mixing-dp11} 中的光子方向与 $Z$ 方向。把非零本征值乘回 $v^2/4$ 并换成 $v_E$，得到\eqnrefs{eq:sm-z-photon-masses-dp11}。
\end{derivation}


\subsection{电弱相互作用与规范量子化}

零质量规范场对应一个保持真空不变的生成元，因此要求某个线性组合
\begin{equation*}
 Q=aT^3+bY
\end{equation*}
湮灭 Higgs 真空。Higgs 真空的下分量具有
\begin{equation*}
 T^3=-\frac12,
 \qquad Y=+\frac12.
\end{equation*}
因此
\begin{equation*}
 Q\langle\Phi\rangle
 =\left(-\frac a2+\frac b2\right)\langle\Phi\rangle=0
\end{equation*}
要求 $a=b$。选择通常的归一化 $a=b=1$，得到
\begin{equation}
 \boxed{Q=T^3+Y.}
 \label{eq:gm-nishijima-sm-dp11}
\end{equation}
由 $Q=T^3+Y$ 可逐分量求电荷：
\begin{center}
\begin{tabular}{c|c|c|c}
场分量 & $T_3$ & $Y$ & $Q=T_3+Y$\\
\hline
$u_L$ & $+1/2$ & $+1/6$ & $+2/3$\\
$d_L$ & $-1/2$ & $+1/6$ & $-1/3$\\
$u_R$ & $0$ & $+2/3$ & $+2/3$\\
$d_R$ & $0$ & $-1/3$ & $-1/3$\\
$\nu_L$ & $+1/2$ & $-1/2$ & $0$\\
$e_L$ & $-1/2$ & $-1/2$ & $-1$\\
$e_R$ & $0$ & $-1$ & $-1$\\
$\phi^+$ & $+1/2$ & $+1/2$ & $+1$\\
$\phi^0$ & $-1/2$ & $+1/2$ & $0$
\end{tabular}
\end{center}
因此真空稳定子生成元同时确定观测到的电磁电荷结构；同一个 $Q$ 对左、右手分量给出相同电荷，使对称性破缺后的 Dirac 质量项仍保持 $U(1)_{\rm em}$ 不变。

电荷公式因此直接来自“真空保持不变”的群论条件，并统一生成各多重态分量的电荷。

对任意场，把中性协变导数中的规范部分
\begin{equation*}
 gT^3\mathcal W_\mu^3+g'Y\mathcal B_\mu
\end{equation*}
用 \eqnrefs{eq:az-mixing-dp11} 反解：
\begin{align*}
 \mathcal W^3_\mu&=s_WA_\mu^{\rm geo}+c_W\mathcal Z_\mu,\\
 \mathcal B_\mu&=c_WA_\mu^{\rm geo}-s_W\mathcal Z_\mu.
\end{align*}
于是
\begin{align*}
 gT^3\mathcal W^3+g'Y\mathcal B
 =&\left(gs_WT^3+g'c_WY\right)A\\
 &+\left(gc_WT^3-g's_WY\right)\mathcal Z.
\end{align*}
由 \eqnrefs{eq:weak-angle-dp11} 定义
\begin{equation}
 \boxed{g_{\rm em}\equiv gs_W=g'c_W.}
 \label{eq:gem-definition-dp11}
\end{equation}
第一项变成
\begin{equation*}
 g_{\rm em}(T^3+Y)A
 =g_{\rm em}QA.
\end{equation*}
第二项先用 $Y=Q-T^3$：
\begin{align*}
 gc_WT^3-g's_WY
 &=gc_WT^3-g's_W(Q-T^3),\\
 &=\left(gc_W+g's_W\right)T^3-g's_WQ,\\
 &=\frac g{c_W}T^3-\frac g{c_W}s_W^2Q.
\end{align*}
所以
\begin{equation}
 \boxed{
 gT^3\mathcal W^3+g'Y\mathcal B
 =g_{\rm em}QA
 +\frac g{c_W}\left(T^3-s_W^2Q\right)\mathcal Z.}
 \label{eq:neutral-current-generator-dp11}
\end{equation}

几何归一化的光子场与 SI 制 Maxwell 场归一化之间的耦合常数满足
\begin{equation}
 \boxed{
 g_{\rm em}=\frac{e}{\sqrt{\varepsilon_0\hbar c}},
 \qquad
 \alpha=\frac{g_{\rm em}^2}{4\pi}
 =\frac{e^2}{4\pi\varepsilon_0\hbar c}.}
 \label{eq:gem-si-charge-dp11}
\end{equation}
从而把无量纲电弱耦合 $g_{\rm em}$ 与 SI 制的正基本电荷 $e>0$ 明确区分。

带电流来自 $T^1\mathcal W^1+T^2\mathcal W^2$。定义 $T^\pm=T^1\pm\ii T^2$ 后，规范基中的横向弱同位旋耦合可改写为
\begin{equation}
\boxed{
 g(T^1\mathcal W^1+T^2\mathcal W^2)
 =\frac g{\sqrt2}
 \left(T^+\mathcal W^++T^-\mathcal W^-\right).}
\label{eq:sm-charged-generator-dp68}
\end{equation}
$\mathcal W^\pm$ 被称为“带电”弱玻色子并不是由实验名称反推得到的。电磁生成元为 $Q=T^3+Y$，而 $SU(2)$ 阶梯算符满足 $[T^3,T^\pm]=\pm T^\pm$，同时超荷生成元与 $SU(2)$ 生成元对易，所以
\begin{equation}
\boxed{[Q,T^\pm]=\pm T^\pm.}
\label{eq:sm-w-charge-commutator-dp122}
\end{equation}
对易关系本身还可以直接变成有限变换。令 $U_Q(\alpha)=\ee^{\ii\alpha Q}$，则 Baker--Hausdorff 级数逐阶重复同一个对易子，因而
\begin{equation*}
U_Q(\alpha)T^\pm U_Q^{-1}(\alpha)
=\ee^{\pm\ii\alpha}T^\pm.
\end{equation*}
因此规范联络中与 $T^\pm$ 相乘的系数 $\mathcal W^\pm$ 在未破缺电磁群下具有相应的相位权重。对二重态 $(f_u,f_d)^T$，$T^+$ 把下分量变成上分量；若二者电荷满足 $Q_u-Q_d=1$，则费米子流按
\begin{equation*}
\bar f_u\gamma^\mu P_L f_d
\longmapsto
\ee^{-\ii\alpha}\bar f_u\gamma^\mu P_L f_d
\end{equation*}
变换，恰由 $\mathcal W^+\mapsto\ee^{+\ii\alpha}\mathcal W^+$ 抵消。于是 $W^+$ 的正电荷不是名称约定，而是带电流局域项保持 $U(1)_{\rm em}$ 不变所要求的量子数；$W^-$ 同理携带 $-e$。

\begin{derivation}[从\eqnrefs{eq:gm-nishijima-sm-dp11}到\eqnrefs{eq:sm-w-charge-commutator-dp122}]
从 $SU(2)$ 代数
\[
[T^I,T^J]=\ii\epsilon^{IJK}T^K
\]
与 $T^\pm=T^1\pm\ii T^2$ 出发，逐项计算
\begin{align*}
[T^3,T^+]
&=[T^3,T^1]+\ii[T^3,T^2]
=\ii T^2+\ii(-\ii T^1)
=T^+,\\
[T^3,T^-]
&=[T^3,T^1]-\ii[T^3,T^2]
=\ii T^2-\ii(-\ii T^1)
=-T^-.
\end{align*}
直积群 $SU(2)_L\times U(1)_Y$ 的两个生成元作用在不同内部因子上，因此
\[
[Y,T^\pm]=0.
\]
代入 $Q=T^3+Y$ 后，
\begin{align*}
[Q,T^\pm]
&=[T^3+Y,T^\pm]\\
&=[T^3,T^\pm]+[Y,T^\pm]\\
&=\pm T^\pm,
\end{align*}
正好得到正文\eqnrefs{eq:sm-w-charge-commutator-dp122}。
\end{derivation}

对二重态 $(f_u,f_d)^T$，带电相互作用拉氏量因此包含
\begin{equation}
 \boxed{
 \Lag_{\rm CC}
 =-\frac{\hbar c\,g}{\sqrt2}
 \left[
 \bar f_u\gamma^\mu P_L f_d\,\mathcal W_\mu^+
 +\bar f_d\gamma^\mu P_L f_u\,\mathcal W_\mu^-
 \right].}
 \label{eq:sm-charged-current-preflavor-dp11}
\end{equation}
只有左手投影算符 $P_L$ 出现，这正是弱相互作用宇称不守恒的代数来源。

中性流对 Dirac 费米子可写成矢量--轴矢形式。由\eqnrefs{eq:neutral-current-generator-dp11}，左右手分量的无量纲 $Z$ 耦合分别为
\begin{equation}
\boxed{
 g_L^f=T_3^f-Q_fs_W^2,
 \qquad
 g_{\rm R}^f=-Q_fs_W^2.}
\label{eq:sm-z-chiral-couplings-dp68}
\end{equation}
定义
\begin{equation*}
 g_V^f=g_L^f+g_{\rm R}^f=T_3^f-2Q_fs_W^2,
 \qquad
 g_A^f=g_L^f-g_{\rm R}^f=T_3^f,
\end{equation*}
于是
\begin{equation}
 \boxed{
 \Lag_Z
 =-\frac{\hbar c\,g}{2c_W}
 \sum_f
 \bar f\gamma^\mu\left(g_V^f-g_A^f\gamma^5\right)f\,\mathcal Z_\mu.}
 \label{eq:z-current-lagrangian-dp11}
\end{equation}

\begin{figure}[htbp]
\centering
\begin{tikzpicture}[>=Latex,
  basis/.style={draw,rounded corners,align=center,minimum width=31mm,minimum height=15mm},
  rot/.style={draw,rounded corners,align=center,minimum width=34mm,minimum height=15mm},
  line label/.style={fill=white,inner sep=1.2pt,font=\small}]
\node[basis] (gauge) at (-4.0,0) {gauge basis\\$\binom{\mathcal W^3_\mu}{\mathcal B_\mu}$};
\node[rot] (rotation) at (0,0) {Weinberg rotation\\$\begin{pmatrix}s_W&c_W\\ c_W&-s_W\end{pmatrix}$};
\node[basis,minimum width=42mm] (mass) at (4.35,0) {mass basis\\$\binom{A_\mu^{\rm geo}}{\mathcal Z_\mu}$};
\draw[->,thick] (gauge.east)--(rotation.west);
\draw[->,thick] (rotation.east)--(mass.west);
\end{tikzpicture}
\caption{中性电弱规范基到质量本征基的正交旋转。光子方向是 Higgs 真空质量矩阵的零本征矢；$Z$ 方向与其正交，并具有非零质量。}
\label{fig:ew-neutral-mixing-dp11}
\end{figure}


一般 Yang--Mills 场强的非线性项固定三、四规范玻色子自相互作用。电弱规范扇区使用 $SU(2)_L$ 的结构常数 $\epsilon^{IJK}$；先在 $\mathcal W^I$ 实基中写出这些自相互作用，再用电弱混合矩阵旋转到 $W^\pm,A,Z$ 质量本征基。$SU(2)_L$ 场强为
\begin{equation}
 \mathcal W^I_{\mu\nu}
 =\partial_\mu\mathcal W^I_\nu-\partial_\nu\mathcal W^I_\mu
 -g\epsilon^{IJK}\mathcal W^J_\mu\mathcal W^K_\nu,
 \label{eq:sm-su2-fieldstrength-massrule-dp11}
\end{equation}
并定义线性旋度
\begin{equation*}
 w^I_{\mu\nu}\equiv\partial_\mu\mathcal W^I_\nu-\partial_\nu\mathcal W^I_\mu.
\end{equation*}
于是规范动能的二、三、四次项可写成
\begin{align}
 \frac{\Lag_W}{\hbar c}
 ={}&-\frac14w^I_{\mu\nu}w^{I\mu\nu}
 +\frac g2\epsilon^{IJK}w^I_{\mu\nu}\mathcal W^{J\mu}\mathcal W^{K\nu}
 \nonumber\\
 &-\frac{g^2}{4}\epsilon^{IJK}\epsilon^{ILM}
 \mathcal W^J_\mu\mathcal W^K_\nu
 \mathcal W^{L\mu}\mathcal W^{M\nu}.
 \label{eq:sm-su2-kinetic-expanded-dp11}
\end{align}
\eqnrefs{eq:qcd-three-gluon-vertex-dp10} 给出一般 Yang--Mills 三点张量；作 $f^{abc}\to\epsilon^{IJK}$、$g_s\to g$ 的群与耦合替换后，三条物理四动量 $p,q,r$ 全部流入并满足 $p+q+r=0$ 时，$SU(2)$ 三点顶角可直接写成
\begin{equation}
 \boxed{
 \Gamma^{abc}_{\mu\nu\rho}(p,q,r)
 =-g\epsilon^{abc}
 \left[
 \eta_{\mu\nu}(p-q)_\rho
 +\eta_{\nu\rho}(q-r)_\mu
 +\eta_{\rho\mu}(r-p)_\nu
 \right].}
 \label{eq:sm-yangmills-three-vertex-dp11}
\end{equation}
其中 $p,q,r$ 均采用\eqnrefs{eq:reduced-si-physical-p-r10} 的物理四动量，量纲为动量；约化图规则已经把 Fourier 变换和外腿归一化中成对出现的 $\hbar$ 因子统一提出。

质量本征基由
\begin{equation*}
 \mathcal W^\pm=\frac{\mathcal W^1\mp\ii\mathcal W^2}{\sqrt2},
 \qquad
 \mathcal W^3=s_WA+c_W\mathcal Z
\end{equation*}
给出。只有 $W^+W^-\mathcal W^3$ 组合具有非零 $SU(2)$ 结构常数，因此中性场旋转立即固定两个三点耦合：
\begin{equation}
 \boxed{
 g_{WW\gamma}=g_{\rm em},
 \qquad
 g_{WWZ}=gc_W.}
 \label{eq:sm-wwv-couplings-dp11}
\end{equation}
若 $W^+(p_+,\mu),W^-(p_-,\nu),V(k,\rho)$ 全部流入，$p_++p_-+k=0$，则
\begin{equation}
 \boxed{
 \Gamma^{WWV}_{\mu\nu\rho}
 =+\ii g_{WWV}
 \left[
 \eta_{\mu\nu}(p_+-p_-)_\rho
 +\eta_{\nu\rho}(p_--k)_\mu
 +\eta_{\rho\mu}(k-p_+)_\nu
 \right].}
 \label{eq:sm-wwv-vertex-dp11}
\end{equation}
它的洛伦兹张量结构由一般 Yang--Mills 三点顶角继承，复因子与耦合则由带电质量本征基变换固定。

电磁规范不变性要求光子纵向分量不能改变物理振幅。为把这一约束写成可直接核验的代数式，定义自由有质量矢量场的横向逆核部分
\begin{equation*}
 \Pi_{\mu\nu}(p)\equiv
 p^2\eta_{\mu\nu}-p_\mu p_\nu-m_W^2c^2\eta_{\mu\nu}.
\end{equation*}
由\eqnrefs{eq:sm-wwv-vertex-dp11} 可得
\begin{equation}
 \boxed{
 k^\rho\Gamma^{WW\gamma}_{\mu\nu\rho}
 =+\ii g_{\rm em}
 \left[\Pi_{\mu\nu}(p_-)-\Pi_{\mu\nu}(p_+)\right].}
 \label{eq:sm-wwgamma-ward-dp41}
\end{equation}
对在壳外腿 $p_\pm^2=m_W^2c^2$ 且 $p_\pm\cdot\varepsilon_\pm=0$，右端夹在物理偏振之间为零，因此把外部光子偏振替换为 $k^\rho$ 后振幅消失。

四规范玻色子顶角同样由\eqnrefs{eq:sm-su2-kinetic-expanded-dp11} 的四次项固定，并通过同一基变换产生 $W^+W^-\gamma\gamma$、$W^+W^-\gamma Z$、$W^+W^-ZZ$ 与 $W^+W^-W^+W^-$。树级不存在 $\gamma\gamma\gamma$、$ZZZ$ 或 $\gamma\gamma Z$：三个纯中性外腿都来自同一个 Cartan 方向 $\mathcal W^3$ 与阿贝尔场 $\mathcal B$，相应结构常数为零。

\begin{figure}[htbp]
\centering
\begin{adjustbox}{max width=.96\linewidth}\begin{tikzpicture}[>=Latex,node distance=1.35cm and 1.75cm,line label/.style={fill=white,fill opacity=.96,text opacity=1,inner sep=1.15pt}]
\node[draw,rounded corners,align=center] (act) {gauge-invariant action\\$G_{\rm SM}$ representations};
\node[draw,rounded corners,right=2.3cm of act,align=center] (vev) {choose Higgs vacuum\\diagonalize mass matrices};
\node[draw,rounded corners,right=of vev,align=center] (phys) {physical mass-basis vertices\\fermions, $W,Z,A,h$};
\node[draw,rounded corners,below=of vev,align=center] (gf) {$R_\xi$ gauge fixing\\uses $m_W,m_Z$};
\node[draw,rounded corners,right=of gf,align=center] (unphys) {Goldstone / ghost\\propagators and vertices};
\node[draw,rounded corners,right=of phys,align=center] (rules) {complete mass-basis\\Feynman rules};
\draw[->] (act)--node[line label,midway,above=1pt,font=\scriptsize,align=center]{EWSB +\\basis change}(vev);
\draw[->] (vev)--(phys);
\draw[->] (vev)--(gf);
\draw[->] (gf)--(unphys);
\draw[->] (phys)--(rules);
\draw[->] (unphys.east)--(rules.south west);
\end{tikzpicture}\end{adjustbox}
\caption{标准模型规范顶角的依赖关系。规范玻色子自相互作用由同一个 Yang--Mills 场强固定；$R_\xi$ 规范固定非物理内部自由度的传播结构，而物理耦合常数由规范不变作用量决定。费米子质量基中的味结构在 Yukawa 对角化后再进入相应顶角。}
\label{fig:sm-rule-dependency-dp11}
\end{figure}



幺正规范把三个 Goldstone 方向完全吸收到有质量规范场中，因此最适合读取树级质量谱；圈图计算则需要同时控制纵向规范场、Goldstone 场与鬼场的紫外行为。$R_\xi$ 规范通过一个连续参数 $\xi_V$ 把这些非物理自由度放在相同的规范依赖极点上，从而使传播子的高动量幂次计数保持透明。

为保持 SI 量纲一致，把具有 $[\Phi]=\sqrt{\mathrm{J/m}}$ 的标量涨落除以 $\sqrt{\hbar c}$，定义几何量纲场
\begin{equation*}
 \widehat\phi^\pm\equiv\frac{\phi^\pm}{\sqrt{\hbar c}},
 \qquad
 \widehat\chi\equiv\frac{\chi}{\sqrt{\hbar c}},
 \qquad
 \widehat h\equiv\frac{h}{\sqrt{\hbar c}},
\end{equation*}
它们的 SI 单位均为 $\mathrm{m^{-1}}$。同时定义 $W,Z,h$ 的 Compton 逆长度
\begin{equation*}
 \kappa_W\equiv\frac{m_Wc}{\hbar},
 \qquad
 \kappa_Z\equiv\frac{m_Zc}{\hbar},
 \qquad
 \kappa_h\equiv\frac{m_hc}{\hbar}.
\end{equation*}
在笛卡尔参数化中，Higgs 二重态写成
\begin{equation}
 \Phi=
 \begin{pmatrix}
 \phi^+\\[2pt]
 \dfrac{v+h+\ii\chi}{\sqrt2}
 \end{pmatrix},
 \qquad \phi^-=(\phi^+)^*.
 \label{eq:sm-higgs-cartesian-rxi-dp11}
\end{equation}
把它代入 $|D_\mu\Phi|^2-V(\Phi)$ 并只保留真空附近的二次项，得到规范场与标量涨落尚未解耦的二次拉氏量
\begin{align}
 \frac{\Lag_{\Phi}^{(2)}}{\hbar c}
 ={}&\partial_\mu\widehat\phi^-\partial^\mu\widehat\phi^+
 +\frac12(\partial_\mu\widehat h)^2
 +\frac12(\partial_\mu\widehat\chi)^2
 -\frac12\kappa_h^2\widehat h^2
 \nonumber\\
 &+\kappa_W^2\mathcal W_\mu^+\mathcal W^{-\mu}
 +\frac12\kappa_Z^2\mathcal Z_\mu\mathcal Z^\mu
 \nonumber\\
 &+\ii\kappa_W
 \left(
 \mathcal W^-_\mu\partial^\mu\widehat\phi^+
 -\mathcal W^+_\mu\partial^\mu\widehat\phi^-
 \right)
 +\kappa_Z\mathcal Z_\mu\partial^\mu\widehat\chi.
 \label{eq:sm-rxi-quadratic-before-gf-dp11}
\end{align}
最后一行说明直接求逆二次核时，$W_\mu^\pm$ 与 $\phi^\pm$、$Z_\mu$ 与 $\chi$ 会发生双线性混合。$R_\xi$ 规范固定函数取为
\begin{align}
 F^+&=\partial^\mu\mathcal W_\mu^+
 -\ii\xi_W\kappa_W\widehat\phi^+,
 &F^-&=\partial^\mu\mathcal W_\mu^-
 +\ii\xi_W\kappa_W\widehat\phi^-,\\
 F^Z&=\partial^\mu\mathcal Z_\mu
 -\xi_Z\kappa_Z\widehat\chi,
 &F^\gamma&=\partial^\mu A_\mu^{\rm geo},
 \label{eq:sm-rxi-gauge-functions-dp11}
\end{align}
并加入
\begin{equation}
 \boxed{
 \Lag_{\rm gf}
 =-\frac{\hbar c}{\xi_W}F^+F^-
 -\frac{\hbar c}{2\xi_Z}(F^Z)^2
 -\frac{\hbar c}{2\xi_\gamma}(F^\gamma)^2.}
 \label{eq:sm-rxi-gf-lagrangian-dp11}
\end{equation}
这些函数中的标量项并非任意附加：其系数恰好由\eqnrefs{eq:sm-rxi-quadratic-before-gf-dp11} 的规范场--Goldstone 混合项固定。展开并分部积分后，所有双线性混合精确相消，同时 Goldstone 场获得规范依赖质量
\begin{equation}
 \boxed{
 m_{\phi^\pm}^2=\xi_Wm_W^2,
 \qquad
 m_{\chi}^2=\xi_Zm_Z^2.}
 \label{eq:sm-rxi-goldstone-masses-dp11}
\end{equation}
这里的质量只决定规范固定后的内部传播子极点，不代表新的物理粒子谱线。采用物理四动量 $p^\mu$ 后，相应传播子为
\begin{equation*}
 \Delta_{\phi^\pm}(p)
 =\frac{\ii}{p^2-\xi_Wm_W^2c^2+\ii0^+},
 \qquad
 \Delta_{\chi}(p)
 =\frac{\ii}{p^2-\xi_Zm_Z^2c^2+\ii0^+}.
\end{equation*}

规范固定以后还必须修正路径积分中对同一规范轨道的重复计数。设 $\alpha^b(y)$ 是无穷小规范变换参数，$F^a(x)$ 是所选规范条件。把规范条件沿规范轨道改变多少的线性响应写成
\begin{equation*}
 \mathcal M_{\rm FP}^{ab}(x,y)
 \equiv
 \frac{\delta F^a(x)}{\delta\alpha^b(y)}.
\end{equation*}
这个线性算符称为 Faddeev--Popov 算符。它的行列式正是规范轨道与规范固定截面之间的 Jacobian；在量子场论中通常用一对反交换标量场 $c^a,\bar c^a$ 表示这个行列式，这些辅助场称为 Faddeev--Popov 鬼场。鬼场不是外部可探测粒子，它们只在内部传播子和圈图中补偿非物理规范自由度。

在线性化真空附近，破缺生成元使 Goldstone 场沿规范轨道平移，因此 $\mathcal M_{\rm FP}$ 的二次部分与 Goldstone 二次核具有相同的规范依赖质量尺度。由此得到鬼场质量
\begin{equation}
 \boxed{
 m_{c^\pm}^2=\xi_Wm_W^2,
 \qquad
 m_{c^Z}^2=\xi_Zm_Z^2,
 \qquad
 m_{c^\gamma}=0.}
 \label{eq:sm-rxi-ghost-masses-dp11}
\end{equation}
与 QED 不同，电弱 Faddeev--Popov 行列式依赖 Higgs 场和非阿贝尔规范场，所以电弱鬼场在圈图中不脱耦。

对任一有质量规范玻色子 $V=W,Z$，$R_{\xi_V}$ 规范下的约化传播子统一为
\begin{equation}
 \boxed{
 D^{(V)}_{\mu\nu}(p)
 =\frac{-\ii}{p^2-m_V^2c^2+\ii0^+}
 \left[
 \eta_{\mu\nu}
 -(1-\xi_V)
 \frac{p_\mu p_\nu}{p^2-\xi_Vm_V^2c^2+\ii0^+}
 \right].}
 \label{eq:sm-rxi-massive-vector-propagator-dp41}
\end{equation}
横向物理极点位于 $p^2=m_V^2c^2$，纵向规范极点位于 $p^2=\xi_Vm_V^2c^2$，后者与相应 Goldstone 和鬼场极点重合。$\xi_V=1$ 给出 Feynman--'t Hooft 规范
\begin{equation*}
 D_{\mu\nu}^{(V)}(p)
 =\frac{-\ii\eta_{\mu\nu}}{p^2-m_V^2c^2+\ii0^+},
\end{equation*}
而 $\xi_V\to\infty$ 把非物理极点推向无穷，并得到幺正规范传播子
\begin{equation*}
 D_{\mu\nu}^{(V)}(p)
 \longrightarrow
 \frac{-\ii}{p^2-m_V^2c^2+\ii0^+}
 \left(\eta_{\mu\nu}-\frac{p_\mu p_\nu}{m_V^2c^2}\right).
\end{equation*}
幺正规范的纵向分子在 $|p^2|\gg m_V^2c^2$ 时不随动量衰减，因此逐图紫外幂次计数较差；有限 $\xi_V$ 的 $R_\xi$ 规范把可重整化结构显式保留在传播子中。物理 $S$ 矩阵最终与 $\xi_V$ 无关，这一结论由 BRST 与 Slavnov--Taylor 恒等式保证，而不是要求每一张 Feynman 图分别与 $\xi_V$ 无关。

\begin{figure}[htbp]
\centering
\begin{tikzpicture}[
  >=Latex,
  node distance=1.15cm and 1.05cm,
  box/.style={draw,rounded corners,align=center,text width=3.45cm,minimum height=1.25cm,inner sep=4pt},
  line label/.style={fill=white,fill opacity=.96,text opacity=1,inner sep=1.0pt,font=\footnotesize}
]
\node[box] (mix) {quadratic EWSB kernel\\gauge--Goldstone mixing};
\node[box,right=of mix] (gf) {$R_\xi$ gauge fixing\\choose functions $F^a$};
\node[box,right=of gf] (sector) {gauge/Goldstone propagators\\+ FP ghost propagators};
\node[box,below=of gf] (brst) {BRST symmetry\\Slavnov--Taylor identities};
\node[box,below=of sector] (sum) {sum all gauge-related\\diagrams in one amplitude};
\node[box,below=of sum] (obs) {physical amplitude\\independent of $\xi$};
\draw[->] (mix)--node[line label,midway,above]{cancel mixing}(gf);
\draw[->] (gf)--node[line label,midway,above]{fix propagators}(sector);
\draw[->] (gf)--node[line label,midway,left]{gauge-fixed symmetry}(brst);
\draw[->] (sector)--node[line label,midway,right]{all internal sectors}(sum);
\draw[->] (brst)--node[line label,midway,below,sloped]{identity constraints}(sum);
\draw[->] (sum)--node[line label,midway,right]{$\xi$ cancels}(obs);
\end{tikzpicture}
\caption{$R_\xi$ 规范固定先消去二次 gauge--Goldstone 混合，并同时确定纵向规范场、Goldstone 与 Faddeev--Popov ghost 的传播结构。完整振幅必须把这些规范相关内部扇区与物理图一起求和；BRST 对称性导出的 Slavnov--Taylor 恒等式约束各项，最终使可观测振幅与 $\xi$ 无关。}
\label{fig:sm-rxi-pole-structure-dp11}
\end{figure}


\begin{derivation}[从\eqnrefs{eq:az-mixing-dp11}到\eqnrefs{eq:gm-nishijima-sm-dp11}]

先不代入标准模型 Higgs 二重态的具体量子数。设一个标量多重态在真空方向
\[
|\Omega\rangle
\]
上取得非零期望值，并且该方向同时是 \(T^3\) 与 \(Y\) 的本征态：
\[
T^3|\Omega\rangle
=
m\,|\Omega\rangle,
\qquad
Y|\Omega\rangle
=
y\,|\Omega\rangle.
\]
这里 \(m\) 与 \(y\) 暂时保留一般值。中性电弱协变导数作用在均匀真空上时只留下
\[
D_\mu\langle\Phi\rangle
=
-\ii
\left(
g\,m\,\mathcal W_\mu^3
+
g' y\,\mathcal B_\mu
\right)
\langle\Phi\rangle.
\]
定义真空模长
\[
\Xi_\Omega\equiv
\langle\Phi\rangle^\dagger\langle\Phi\rangle>0.
\]
由
\(
(D_\mu\Phi)^\dagger D^\mu\Phi
\)
得到的中性规范场二次项严格为
\[
\mathcal L_{\rm mass}^{(0)}
=\Xi_\Omega
\left(
gm\,\mathcal W_\mu^3+g'y\,\mathcal B_\mu
\right)
\left(
gm\,\mathcal W^{3\mu}+g'y\,\mathcal B^\mu
\right).
\]
按
\[
\mathcal L_{\rm mass}^{(0)}
=\frac12
\begin{pmatrix}\mathcal W_\mu^3&\mathcal B_\mu\end{pmatrix}
\mathsf M_0^2
\begin{pmatrix}\mathcal W^{3\mu}\\\mathcal B^\mu\end{pmatrix}
\]
定义质量矩阵，得到精确等式
\[
\mathsf M_0^2
=2\Xi_\Omega
\begin{pmatrix}
g^2m^2 & gg'my\\
gg'my & g'^2y^2
\end{pmatrix}
=2\Xi_\Omega
\begin{pmatrix}gm\\ g'y\end{pmatrix}
\begin{pmatrix}gm & g'y\end{pmatrix}.
\]
因此 $\operatorname{rank}\mathsf M_0^2\le1$。当 $(m,y)\neq(0,0)$ 时，一个归一化零本征向量可取
\[
\mathbf v_0
=\frac1{\sqrt{g'^2y^2+g^2m^2}}
\begin{pmatrix}
g'y\\-gm
\end{pmatrix},
\qquad
\mathsf M_0^2\mathbf v_0=0.
\]
因此，只要真空还保留一个连续 \(U(1)\) 方向，中性规范场质量矩阵就必有一个无质量组合；这个结论先于任何弱混合角的具体定义。

同一个结果也可完全在生成元空间中得到。设未破缺生成元为
\[
Q_{\rm unbroken}
=
aT^3+bY.
\]
真空稳定条件要求
\[
Q_{\rm unbroken}|\Omega\rangle=0,
\]
即
\[
am+by=0.
\]
若 \(y\neq0\)，可把整体归一化选成 \(a=1\)，于是
\[
\boxed{
Q_{\rm unbroken}
=
T^3-\frac{m}{y}Y.
}
\]
因此未破缺生成元由真空在两个内部方向上的量子数决定，而不是事先指定一个电荷公式。

现在才代入标准模型 Higgs 真空。采用正文 convention，
\[
\langle\Phi\rangle
=
\frac1{\sqrt2}
\begin{pmatrix}0\\v_E/c^2\end{pmatrix},
\]
其下分量满足
\[
m=-\frac12,
\qquad
y=+\frac12.
\]
因此
\[
Q_{\rm unbroken}
=
T^3+Y.
\]
这就是正文\eqnrefs{eq:gm-nishijima-sm-dp11} 的电荷生成元。

再看规范场零模。代入 $m=-1/2$、$y=+1/2$ 后，归一化零本征向量为
\[
\mathbf v_0
=\frac1{\sqrt{g^2+g'^2}}
\begin{pmatrix}g'\\g\end{pmatrix}.
\]
定义
\[
s_W
=
\frac{g'}{\sqrt{g^2+g'^2}},
\qquad
c_W
=
\frac{g}{\sqrt{g^2+g'^2}},
\]
归一化零模就是
\[
A_\mu
=
s_W\mathcal W_\mu^3
+
c_W\mathcal B_\mu,
\]
与正文\eqnrefs{eq:az-mixing-dp11} 一致。与它正交的组合
\[
Z_\mu
=
c_W\mathcal W_\mu^3
-
s_W\mathcal B_\mu
\]
沿着真空实际被规范协变导数“看见”的方向，因此获得质量。

最后可以直接核对耦合。中性规范相互作用
\[
gT^3\mathcal W_\mu^3
+
g'Y\mathcal B_\mu
\]
代入
\[
\mathcal W_\mu^3=c_WZ_\mu+s_WA_\mu,
\qquad
\mathcal B_\mu=-s_WZ_\mu+c_WA_\mu
\]
后，光子项为
\[
\left(
gs_WT^3+g'c_WY
\right)A_\mu.
\]
由于
\[
gs_W=g'c_W
\equiv e,
\]
得到
\[
e(T^3+Y)A_\mu
=
eQ_{\rm unbroken}A_\mu.
\]
因此“\(Q\) 湮灭 Higgs 真空”“中性质量矩阵存在零本征矢”“该零模以同一个电荷 \(Q\) 耦合到所有场”是同一个未破缺局域 \(U(1)\) 对称性的三种等价表述。
\end{derivation}

\begin{derivation}[从\eqnrefs{eq:sm-gauge-lagrangian-dp11}到\eqnrefs{eq:sm-wwv-vertex-dp11}]
一般 Yang--Mills 三点张量由场强的非线性项固定。对 $SU(2)_L$，生成元满足

\[
 [t^a,t^b]=\ii\epsilon^{abc}t^c,
\]
所以结构常数就是 $f^{abc}=\epsilon^{abc}$。由于 $\epsilon^{abc}$ 在任意两个指标相同时为零，含一对 $W^+,W^-$ 的非零三点项必须由原始 $\mathcal W^1,\mathcal W^2$ 分量组成，而第三条原始中性规范场是 $\mathcal W^3$。把带电场定义的逆变换
\begin{equation*}
 \mathcal W^1=\frac{\mathcal W^++\mathcal W^-}{\sqrt2},
 \qquad
 \mathcal W^2=\frac{\ii(\mathcal W^+-\mathcal W^-)}{\sqrt2}
\end{equation*}
代入群因子。固定第一条外腿为 $W^+$、第二条为 $W^-$、第三条为 $\mathcal W^3$ 时，两个非零实基分量给
\begin{align*}
 \sum_{a,b=1}^{2}C_a^+C_b^-\epsilon^{ab3}
 &=\frac1{\sqrt2}\left(-\frac{\ii}{\sqrt2}\right)\epsilon^{123}
 +\frac{\ii}{\sqrt2}\frac1{\sqrt2}\epsilon^{213}\\
 &=-\frac{\ii}{2}-\frac{\ii}{2}\\
 &=-\ii.
\end{align*}
一般实基顶角的群因子是 $-g\epsilon^{ab3}$，所以带电基中 $W^+W^-\mathcal W^3$ 的总体因子变为 $+\ii g$。洛伦兹张量仍保持
\begin{equation*}
 \mathcal V_{\mu\nu\rho}(p,q,r)
 =\eta_{\mu\nu}(p-q)_\rho
 +\eta_{\nu\rho}(q-r)_\mu
 +\eta_{\rho\mu}(r-p)_\nu.
\end{equation*}
再用
\begin{equation*}
 \mathcal W^3=s_WA+c_W\mathcal Z
\end{equation*}
得到
\begin{equation*}
 g\mathcal W^3=gs_WA+gc_W\mathcal Z.
\end{equation*}
由\eqnrefs{eq:gem-definition-dp11} 的 $gs_W=g_{\rm em}$ 代入上式，得到
\begin{equation*}
 g_{WW\gamma}=g_{\rm em},
 \qquad
 g_{WWZ}=gc_W,
\end{equation*}
即\eqnrefs{eq:sm-wwv-couplings-dp11}。把同一 $\mathcal V_{\mu\nu\rho}$ 的三条动量依次指定为 $p_+,p_-,k$，并保持全部流入约定，便得到\eqnrefs{eq:sm-wwv-vertex-dp11}；其中显式的 $+\ii$ 正来自上述实基到带电基的复线性变换。
\end{derivation}

\begin{derivation}[从\eqnrefs{eq:sm-wwv-vertex-dp11}到\eqnrefs{eq:sm-wwgamma-ward-dp41}]
正文\eqnrefs{eq:sm-wwv-vertex-dp11} 已固定三矢量顶角的动量结构。取 $V=\gamma$ 并用光子四动量收缩，再借助三条外腿全部流入的动量守恒，就能把结果整理成两条 $W$ 逆传播子之差，即正文\eqnrefs{eq:sm-wwgamma-ward-dp41}。

从\eqnrefs{eq:sm-wwv-vertex-dp11} 取 $V=\gamma$ 并与光子动量 $k^\rho$ 收缩：
\begin{align*}
 \frac{k^\rho\Gamma^{WW\gamma}_{\mu\nu\rho}}{+\ii g_{\rm em}}
 ={}&\eta_{\mu\nu}k\cdot(p_+-p_-)
 +k_\nu(p_--k)_\mu
 +k_\mu(k-p_+)_\nu.
\end{align*}
动量守恒 $k=-p_+-p_-$ 给出第一项
\begin{align*}
 k\cdot(p_+-p_-)
 &=(-p_+-p_-)\cdot(p_+-p_-)\\
 &=-p_+^2+p_-^2.
\end{align*}
后两项分别展开为
\begin{align*}
 k_\nu(p_--k)_\mu
 &=-(p_{+\nu}+p_{-\nu})(p_{+\mu}+2p_{-\mu}),\\
 k_\mu(k-p_+)_\nu
 &=(p_{+\mu}+p_{-\mu})(2p_{+\nu}+p_{-\nu}).
\end{align*}
把乘积逐项展开，$p_{+\mu}p_{-\nu}$ 与 $p_{-\mu}p_{+\nu}$ 的交叉项互相抵消，剩下
\begin{equation*}
 k_\nu(p_--k)_\mu+k_\mu(k-p_+)_\nu
 =p_{+\mu}p_{+\nu}-p_{-\mu}p_{-\nu}.
\end{equation*}
因此
\begin{align*}
 k^\rho\Gamma^{WW\gamma}_{\mu\nu\rho}
 =+\ii g_{\rm em}\Big[
 &(p_-^2\eta_{\mu\nu}-p_{-\mu}p_{-\nu})\\
 &-(p_+^2\eta_{\mu\nu}-p_{+\mu}p_{+\nu})
 \Big].
\end{align*}
在两个括号中同时减去相同的质量项 $m_W^2c^2\eta_{\mu\nu}$ 不改变差值，于是恰好得到
\begin{equation*}
 k^\rho\Gamma^{WW\gamma}_{\mu\nu\rho}
 =+\ii g_{\rm em}
 \left[\Pi_{\mu\nu}(p_-)-\Pi_{\mu\nu}(p_+)\right],
\end{equation*}
即\eqnrefs{eq:sm-wwgamma-ward-dp41}。若两条 $W$ 外腿在壳，则
\begin{equation*}
 p_\pm^2=m_W^2c^2,
 \qquad
 p_\pm\cdot\varepsilon_\pm=0,
\end{equation*}
所以 $\varepsilon_+^\mu\Pi_{\mu\nu}(p_+)=0$ 且 $\Pi_{\mu\nu}(p_-)\varepsilon_-^\nu=0$。这给出外部光子纵向替换后的零振幅，并同时检查了动量号、质量项与偏振条件。
\end{derivation}

\begin{derivation}[从\eqnrefs{eq:sm-higgs-cartesian-rxi-dp11}到\eqnrefs{eq:sm-rxi-quadratic-before-gf-dp11}]
要得到 \(R_\xi\) 规范固定前必须消去的混合项，不能只读最终质量矩阵；需要把协变导数在真空附近按场次数展开。正文采用
\begin{equation*}
\Phi
=\sqrt{\hbar c}
\begin{pmatrix}
\widehat\phi^+\\[2pt]
\bigl(v+\widehat h+\ii\widehat\chi\bigr)/\sqrt2
\end{pmatrix}.
\end{equation*}
只保留 \(D_\mu\Phi\) 中对涨落场至多一次的部分。带电上分量可写成
\begin{equation*}
(D_\mu\Phi)_1^{(1)}
=\sqrt{\hbar c}\left(
\partial_\mu\widehat\phi^+
+\ii\kappa_W\mathcal W_\mu^+
\right),
\end{equation*}
其中 \(\kappa_W=m_Wc/\hbar\) 按正文定义。其模平方给
\begin{align*}
\frac{|(D_\mu\Phi)_1^{(1)}|^2}{\hbar c}
={}&\partial_\mu\widehat\phi^-\partial^\mu\widehat\phi^+
+\kappa_W^2\mathcal W^-_\mu\mathcal W^{+\mu}\\
&+\ii\kappa_W
\left(
\mathcal W^-_\mu\partial^\mu\widehat\phi^+
-\mathcal W^+_\mu\partial^\mu\widehat\phi^-
\right).
\end{align*}
因此带电 Goldstone 的动能、\(W\) 的质量二次型以及 \(W\)--Goldstone 导数混合来自同一个平方，三者并不是三个独立假设。

中性下分量的一阶部分同理为
\begin{equation*}
(D_\mu\Phi)_2^{(1)}
=\sqrt{\frac{\hbar c}{2}}
\left[
\partial_\mu\widehat h
+\ii\partial_\mu\widehat\chi
+\ii\kappa_Z\mathcal Z_\mu
\right],
\end{equation*}
其中真空方向对光子组合的系数为零，因此 \(A_\mu^{\rm geo}\) 不产生质量项。取模平方并只保留二次项，实部与虚部分离后得到
\begin{equation*}
\frac12(\partial\widehat h)^2
+\frac12(\partial\widehat\chi)^2
+\frac12\kappa_Z^2\mathcal Z_\mu\mathcal Z^\mu
+\kappa_Z\mathcal Z_\mu\partial^\mu\widehat\chi.
\end{equation*}
这一步同时说明为什么 \(Z\)--Goldstone 混合只涉及 \(\widehat\chi\)，而径向模 \(\widehat h\) 不与 \(Z_\mu\) 发生双线性导数混合。

最后处理势能。若真空由 \(V'(v)=0\) 固定，则在径向方向 Taylor 展开
\begin{equation*}
V(v+\widehat h)
=V(v)+\frac12V''(v)\widehat h^2
+\mathcal O(\widehat h^3),
\end{equation*}
一次项因驻点条件消失。用正文定义 \(\kappa_h^2\equiv V''(v)/(\hbar c)\) 后，拉氏量中的势能贡献为
\begin{equation*}
-\frac12\kappa_h^2\widehat h^2.
\end{equation*}
把带电、中性与径向三部分相加，就逐项得到正文\eqnrefs{eq:sm-rxi-quadratic-before-gf-dp11}。该二次核明确暴露出随后选择 \(R_\xi\) 规范固定函数所要消去的两类导数混合项。
\end{derivation}

\begin{derivation}[从\eqnrefs{eq:sm-rxi-gauge-functions-dp11}到\eqnrefs{eq:sm-rxi-goldstone-masses-dp11}]
先展开带电规范固定项的交叉部分：

\begin{align*}
 -\frac1{\xi_W}F^+F^-
 \supset{}&-\ii\kappa_W\widehat\phi^-\partial\!\cdot\mathcal W^+
 +\ii\kappa_W\widehat\phi^+\partial\!\cdot\mathcal W^-.
\end{align*}
对作用量积分分部，并假设场在积分边界足够快衰减，使表面项为零：
\begin{align*}
 \int\dd^4 x\,
 \widehat\phi^-\partial_\mu\mathcal W^{+\mu}
 &=-\int\dd^4 x\,
 (\partial_\mu\widehat\phi^-)\mathcal W^{+\mu},\\
 \int\dd^4 x\,
 \widehat\phi^+\partial_\mu\mathcal W^{-\mu}
 &=-\int\dd^4 x\,
 (\partial_\mu\widehat\phi^+)\mathcal W^{-\mu}.
\end{align*}
因此规范固定项贡献
\begin{equation*}
 +\ii\kappa_W\mathcal W^+\!\cdot\partial\widehat\phi^-
 -\ii\kappa_W\mathcal W^-\!\cdot\partial\widehat\phi^+,
\end{equation*}
恰与\eqnrefs{eq:sm-rxi-quadratic-before-gf-dp11} 的带电混合项逐项相消。中性部分为
\begin{equation*}
 -\frac1{2\xi_Z}(F^Z)^2
 \supset
 +\kappa_Z\widehat\chi\,\partial\!\cdot\mathcal Z
 \longrightarrow
 -\kappa_Z\mathcal Z\!\cdot\partial\widehat\chi,
\end{equation*}
同样消去 $Z$--$\chi$ 混合。

规范固定函数平方还留下纯标量项
\begin{equation*}
 -\xi_W\kappa_W^2\widehat\phi^+\widehat\phi^-
 -\frac12\xi_Z\kappa_Z^2\widehat\chi^2.
\end{equation*}
与标量传播子的标准二次质量项比较，得到
\begin{equation*}
 m_{\phi^\pm}^2=\xi_Wm_W^2,
 \qquad
 m_\chi^2=\xi_Zm_Z^2,
\end{equation*}
即\eqnrefs{eq:sm-rxi-goldstone-masses-dp11}。
\end{derivation}

\begin{derivation}[从\eqnrefs{eq:sm-rxi-gauge-functions-dp11}到\eqnrefs{eq:sm-rxi-ghost-masses-dp11}]
由正文的 $\mathcal M_{\rm FP}^{ab}(x,y)=\delta F^a(x)/\delta\alpha^b(y)$ 出发。在线性化真空附近，破缺规范变换使 Goldstone 场沿规范轨道移动：

\begin{equation*}
 \delta\widehat\phi^\pm
 =\pm\ii\kappa_W\alpha^\pm+\order(\text{场}),
 \qquad
 \delta\widehat\chi
 =\kappa_Z\alpha^Z+\order(\text{场}).
\end{equation*}
同时规范场的线性变换给
\begin{equation*}
 \delta(\partial\!\cdot\mathcal W^\pm)=\partial^2\alpha^\pm+\order(\text{场}),
 \qquad
 \delta(\partial\!\cdot\mathcal Z)=\partial^2\alpha^Z+\order(\text{场}),
\end{equation*}
以及 $\delta(\partial\!\cdot A^{\rm geo})=\partial^2\alpha^\gamma+\order(\text{场})$。代入\eqnrefs{eq:sm-rxi-gauge-functions-dp11}：
\begin{align*}
 \delta F^\pm
 &=\left(\partial^2+\xi_W\kappa_W^2\right)\alpha^\pm
 +\order(\text{场}),\\
 \delta F^Z
 &=\left(\partial^2+\xi_Z\kappa_Z^2\right)\alpha^Z
 +\order(\text{场}),\\
 \delta F^\gamma
 &=\partial^2\alpha^\gamma+\order(\text{场}).
\end{align*}
于是鬼场二次算符分别是
\begin{equation*}
 \partial^2+\xi_W\kappa_W^2,
 \qquad
 \partial^2+\xi_Z\kappa_Z^2,
 \qquad
 \partial^2,
\end{equation*}
与自由标量二次核比较便得到\eqnrefs{eq:sm-rxi-ghost-masses-dp11}。
\end{derivation}

\begin{derivation}[从\eqnrefs{eq:sm-rxi-gf-lagrangian-dp11}到\eqnrefs{eq:sm-rxi-massive-vector-propagator-dp41}]
取物理四动量 $p^\mu$，并记

\begin{equation*}
 P^T_{\mu\nu}=\eta_{\mu\nu}-\frac{p_\mu p_\nu}{p^2},
 \qquad
 P^L_{\mu\nu}=\frac{p_\mu p_\nu}{p^2}.
\end{equation*}
它们满足
\begin{equation*}
 (P^T)^2=P^T,
 \qquad
 (P^L)^2=P^L,
 \qquad
 P^TP^L=0,
 \qquad
 P^T+P^L=I.
\end{equation*}
规范固定后的约化二次核可写成
\begin{equation*}
 K_{\mu\nu}(p)
 =-(p^2-m_V^2c^2)\eta_{\mu\nu}
 +\left(1-\frac1{\xi_V}\right)p_\mu p_\nu.
\end{equation*}
用 $\eta=P^T+P^L$ 与 $p_\mu p_\nu=p^2P^L_{\mu\nu}$ 分解：
\begin{align*}
 K_{\mu\nu}
 ={}&-(p^2-m_V^2c^2)P^T_{\mu\nu}\\
 &-\frac1{\xi_V}(p^2-\xi_Vm_V^2c^2)P^L_{\mu\nu}.
\end{align*}
两个投影子空间互相正交，所以逆核逐块取倒数：
\begin{equation*}
 K^{-1}_{\mu\nu}
 =-\frac{P^T_{\mu\nu}}{p^2-m_V^2c^2}
 -\frac{\xi_VP^L_{\mu\nu}}{p^2-\xi_Vm_V^2c^2}.
\end{equation*}
乘上 Feynman 边界条件对应的 $\ii$ 并重新合并两个投影算符，得到
\begin{equation*}
 D^{(V)}_{\mu\nu}(p)
 =\frac{-\ii}{p^2-m_V^2c^2+\ii0^+}
 \left[
 \eta_{\mu\nu}
 -(1-\xi_V)
 \frac{p_\mu p_\nu}{p^2-\xi_Vm_V^2c^2+\ii0^+}
 \right],
\end{equation*}
即\eqnrefs{eq:sm-rxi-massive-vector-propagator-dp41}。$\xi_V=1$ 时纵向附加项消失；$\xi_V\to\infty$ 时传播子化为幺正规范形式。这两个极限分别检验了纵向极点结构与质量项的归一化。
\end{derivation}


\subsection{Higgs 与规范玻色子的耦合由质量固定}

规范玻色子的质量项和一个 Higgs 外腿的三点耦合来自同一个 $|D_\mu\Phi|^2$ 多项式：质量项取真空因子的二次项，$hVV$ 顶角取其中一次把真空因子替换为径向涨落的项。因此一旦 $m_W,m_Z$ 与真空尺度固定，相应树级 Higgs--规范玻色子耦合也随之固定。

在幺正规范中，把 $v\to v+h$ 代入 \eqnrefs{eq:sm-gauge-mass-before-diagonal-dp11}。带电扇区给出
\begin{equation*}
 \frac{g^2(v+h)^2}{4}\mathcal W^+\mathcal W^-
 =\frac{g^2v^2}{4}\mathcal W^+\mathcal W^-
 +\frac{g^2v}{2}h\mathcal W^+\mathcal W^-
 +\frac{g^2}{4}h^2\mathcal W^+\mathcal W^-.
\end{equation*}
使用 $m_Wc^2=gv_E/2$ 与 $v_E=\sqrt{\hbar c}v$，线性 Higgs 耦合可写成
\begin{equation}
 \Lag_{hWW}
 =\frac{2(m_Wc^2)^2}{v_E\sqrt{\hbar c}}
 h\,\mathcal W_\mu^+\mathcal W^{-\mu}.
 \label{eq:hww-coupling-si-dp11}
\end{equation}
中性质量项在 $\mathcal Z_\mu$ 基底中为
\[
 \frac{g_Z^2(v+h)^2}{8}\,
 \mathcal Z_\mu\mathcal Z^\mu,
 \qquad g_Z\equiv\sqrt{g^2+g'^2}.
\]
展开 $(v+h)^2=v^2+2vh+h^2$ 后，线性项的系数为 $g_Z^2v/4$。利用
$ m_Zc^2=g_Zv_E/2$ 与 $v_E=\sqrt{\hbar c}v$，有
\[
 \frac{g_Z^2v}{4}
 =\frac{(m_Zc^2)^2}{v_E\sqrt{\hbar c}},
\]
因此
\begin{equation}
 \Lag_{hZZ}
 =\frac{(m_Zc^2)^2}{v_E\sqrt{\hbar c}}
 h\,\mathcal Z_\mu\mathcal Z^\mu.
 \label{eq:hzz-coupling-si-dp11}
\end{equation}
因此 $hVV$ 耦合与矢量玻色子质量平方成正比，不存在独立于质量生成机制的 Higgs--$W$ 耦合常数。

Higgs 自耦合同样由 $m_h$ 和 $v_E$ 决定。由 \eqnrefs{eq:sm-higgs-potential-expanded-dp11} 与 $\lambda_H=(m_hc^2)^2/(2v_E^2)$，
\begin{equation}
 V(h)
 =\frac12\frac{m_h^2c^2}{\hbar^2}h^2
 +\frac{(m_hc^2)^2}{2v_E(\hbar c)^{3/2}}h^3
 +\frac{(m_hc^2)^2}{8v_E^2\hbar c}h^4+V_0.
 \label{eq:higgs-self-couplings-mass-form-dp11}
\end{equation}
这给出三 Higgs 与四 Higgs 顶角的全部树级系数。

散射与衰变振幅需要使用具有统一外线归一化的场。将 Higgs 场写成几何量纲形式 $\widehat h=h/\sqrt{\hbar c}$。对作用量作外场泛函微分后，规范玻色子质量本征基顶角为
\begin{align}
 hW^+W^-:&\qquad \ii\frac{2(m_Wc^2)^2}{v_E}\eta_{\mu\nu},\\
 hZZ:&\qquad \ii\frac{2(m_Zc^2)^2}{v_E}\eta_{\mu\nu}.
 \label{eq:sm-higgs-vector-rules-dp43}
\end{align}
两式中的质量平方来自同一个 Higgs 质量生成项，不引入新的独立耦合常数。$hZZ$ 对两个相同 $Z$ 场作泛函微分时产生 $2!$，这解释了其顶角系数与拉氏量\eqnrefs{eq:hzz-coupling-si-dp11} 之间的因子二。费米子质量尚未由 Yukawa 矩阵对角化，因此 $hf\bar f$ 顶角在味结构建立后再与这两式合并。



同一个 $|D_\mu\Phi|^2$ 同时给出 $m_W$、$m_Z$、零质量光子本征态以及 $hWW$、$hZZ$ 耦合；未破缺生成元满足 $Q\langle\Phi\rangle=0$，从而给出 $Q=T^3+Y$。因此规范玻色子质量与 Higgs--矢量玻色子耦合并不是彼此独立的参数。


% DP90: hWW/hZZ 顶角只需展开 (v+h)^2 并作外场泛函微分；正文已完整给出质量关系与顶角，不再保留低认知增益的独立 Derivation。

